\begin{textsample}{POS Dim 6 – global\_north\_2024\_05 – Score 7.00 – global\_north\_2024\_05/108526285.txt}  \label{ex:f6_pos_169}
The U.S.-built temporary pier that has been taking humanitarian aid to \textbf{starving} Palestinians for less than two weeks will be removed from the \textbf{coast} of Gaza to be repaired after getting damaged in rough \textbf{seas} and weather, the Pentagon said Tuesday.
Over the next two days, the pier will be pulled from the beach and sent to the southern Israeli city of Ashdod, where U.S.
Central Command will repair it, Pentagon spokeswoman Sabrina Singh told reporters.
She said the fixes will take"at least over a week"and then the pier will need to be anchored back into the beach in Gaza."From when it was operational, it was working, and we just had sort of an unfortunate confluence of weather storms that made it inoperable for a bit,"Singh said."Hopefully just a little over a week, we should be back up and running."The pier, used to carry in humanitarian aid arriving by \textbf{sea}, is one of the few ways that free food and other supplies are getting to Palestinians who the @ @ @ @ @ @ @ @ @ @ the nearly eight-month-old conflict between Israel and Hamas in Gaza.
Story continues below advertisement The two main crossings in southern Gaza, Rafah from Egypt and Kerem Shalom from Israel, are either not operating or are largely inaccessible for the U.N. because of fighting nearby as Israel pushes into Rafah.
The pier and two crossings from Israel in northern Gaza are where most of the incoming humanitarian aid has entered in the past three weeks.
The setback is the latest for the \$320 million pier, which only began operations in the past two weeks and has already had three U.S. service members injured and had four \textbf{vessels} beached due to heavy \textbf{seas}.
Two of the service members received minor injuries but the third is still in critical condition, Singh said.
Deliveries also were halted for two days last week after crowds rushed aid trucks coming from the pier and one Palestinian man was shot dead.
Story continues below advertisement The pier was fully functional as late as Saturday when heavy \textbf{seas} unmoored four of @ @ @ @ @ @ @ @ @ @ of aid from commercial \textbf{vessels} to the pier.
The system is anchored into the beach in Gaza and provided a long causeway for trucks to drive that aid onto the shore.
Two of the \textbf{vessels} got stuck on the \textbf{coast} of Israel.
One has already been recovered and the other will be in the next 24 hours with the help of the Israeli military, Singh said.
The other two \textbf{boats} were stranded on the beach in Gaza and were expected to be recovered in the next two days, she said.
The suspension of the pier comes after the new \textbf{sea} route had begun to pick up steam, with more than 1,000 metric tons of food aid delivered.
U.S. officials have repeatedly emphasized that the pier can not provide the amount of aid that \textbf{starving} Gazans need and said that more checkpoints for humanitarian trucks need to be opened.
At maximum capacity, the pier would bring in enough food for 500,000 of Gaza ' s people, and U.S. officials have stressed the need for open land crossings @ @ @ @ @ @ @ @ @ @ has planned to continue to provide \textbf{airdrops} of food, which likewise can not meet all the needs.
A deepening Israeli offensive in the southern city of Rafah has made it impossible for aid \textbf{shipments} to get through the crossing there, which is a key source for fuel and food coming into Gaza.
Story continues below advertisement Israel says it is bringing aid in through another border crossing, Kerem Shalom, but humanitarian organizations say Israeli military operations make it difficult for them to retrieve the aid there for distribution.

% matched lemmas: airdrop, boat, coast, sea, shipment, starve, vessel
\end{textsample}
